% ==================================================================================================
% main.tex: the root file. This is the file you compile. This file only decides the ORDER of
%   things. All settings live elsewhere:
% config.tex: your details (title, name, supervisor, ...) and a few layout values (margins, line
%   spacing, link colour)
% preamble/*.tex: package setup, one topic per file
% frontmatter/*.tex: cover, acknowledgements, abstract, contents/lists
% chapters/*.tex: one file per chapter
% backmatter/*.tex: reference list, appendices
% glossary/entries.tex: abbreviations and glossary terms
% references.bib: bibliography database
%
% \input{file}: pastes file.tex in at that point (use for small pieces and preamble/*).
% \include{file}: does the same but always starts a new page. Use it for chapters and appendices
%   only. It cannot be nested.
% ==================================================================================================

% 12pt: body text size
% a4paper: paper size
% oneside: single-sided PDF, identical left/right margins, no blank pages
% openany: chapters start on the next page
% book: provides \frontmatter, \mainmatter and \appendix, which give the roman-then-arabic page
% numbering and chapter handling.
\documentclass[12pt,a4paper,oneside,openany]{book}

% =========================================== preamble =============================================
% =============================================================================
% config.tex: everything specific to you, in one place.
%
% Each line defines a command (a "macro") that the rest of the template uses.
% \newcommand{\name}{text} means: wherever \name appears, print "text".
% Edit only the text inside the SECOND pair of braces; leave the \name alone,
% because other files refer to it.
%
% Where each value is used:
%   \thesistitle       PDF metadata (title shown in the PDF viewer)
%   \covertitle        cover page title
%   \thesisauthor      cover page and PDF metadata
%   \thesisdegree      cover page and PDF metadata
%   \thesisdepartment  cover page
%   \thesissupervisor  cover page
%   \thesisdate        cover page
%   \thesismargin      preamble/layout.tex (page margins)
%   \thesislinespacing preamble/layout.tex (body line spacing)
%   \thesislinkcolour  preamble/links.tex (colour of clickable links)
% =============================================================================

% --- Title -------------------------------------------------------------------
% Plain title with no formatting commands and no \\ line breaks: PDF metadata
% can only hold plain text. This is a working title; change it freely.
\newcommand{\thesistitle}{Detecting Rug-Pull Behaviour in MCP Servers with Windowed REV2 Trust Scores}

% Title as printed on the cover. By default it is the same text, wrapped
% automatically to 2-3 lines. To choose the line breaks yourself, replace
% \thesistitle with the title typed out and \\ where each line should end, e.g.
%   \newcommand{\covertitle}{Detecting Rug-Pull Behaviour\\ in MCP Servers ...}
\newcommand{\covertitle}{\thesistitle}

% --- People and programme ------------------------------------------------------
\newcommand{\thesisauthor}{Paritosh Panda}
\newcommand{\thesisdegree}{MSc Cyber Security Management}
\newcommand{\thesisdepartment}{Warwick Manufacturing Group (WMG)}
\newcommand{\thesissupervisor}{Muhammad Ajmal Azad}
\newcommand{\thesisdate}{September 2026}   % submission month and year

% --- Layout ---------------------------------------------------------------------
% Font size is not here: it is the "12pt" option on the \documentclass line in
% main.tex (allowed values: 10pt, 11pt, 12pt).

% Page margins on all four sides. Any LaTeX length works: 2.5cm, 25mm, 1in.
\newcommand{\thesismargin}{2.5cm}

% Body line spacing, using the setspace package's commands:
%   \singlespacing     single
%   \onehalfspacing    1.5 lines (matches "1.5" in Word)
%   \doublespacing     double
%   \setstretch{1.3}   any other factor
% Footnotes, captions, the cover and the contents/lists stay single-spaced.
\newcommand{\thesislinespacing}{\onehalfspacing}

% --- Links ----------------------------------------------------------------------
% Colour of citations, cross-references, glossary links and URLs. "linkblue"
% is defined in preamble/packages.tex. Use "black" for a printed copy (links
% still work in the PDF, they just aren't coloured). Any xcolor name also
% works (e.g. "blue", "NavyBlue").
\newcommand{\thesislinkcolour}{linkblue}
               % details and layout values
% =============================================================================
% preamble/packages.tex: general-purpose packages used across the document.
%
% \usepackage[options]{name} loads a package (a library of commands).
% Load order occasionally matters; where it does, a comment says why.
% Topic-specific packages live in the other preamble files:
%   fonts.tex (fonts), layout.tex (geometry, setspace, fancyhdr, titlesec),
%   floats.tex (caption, booktabs, algorithm), references.tex (biblatex),
%   links.tex (hyperref, cleveref), glossary.tex (glossaries-extra).
% Everything here works with pdfLaTeX (Overleaf's default), XeLaTeX and
% LuaLaTeX; only fonts.tex behaves differently per compiler.
% =============================================================================

% --- Language and quotation marks ---------------------------------------------------
\usepackage[british]{babel}  % British English: hyphenation patterns, and the
                             % names "Contents", "Chapter", "Appendix", dates
\usepackage{csquotes}        % language-aware quotation marks; biblatex needs it.
                             % \enquote{text} gives 'text' in British style.
\usepackage{etoolbox}        % programming helpers; \appto is used in layout.tex

% --- Maths -------------------------------------------------------------------------
\usepackage{amsmath}      % equation environments (align, gather, cases, ...)
\usepackage{amssymb}      % extra symbols (\mathbb, \leqslant, ...)
\usepackage{mathtools}    % fixes and extensions to amsmath (\coloneqq, ...)

% --- Graphics and colour -----------------------------------------------------------
\usepackage{graphicx}                 % \includegraphics for images
\graphicspath{{figures/}{assets/}}    % folders searched for images, so you can
                                      % write \includegraphics{name} with no path
                                      % and no file extension
\usepackage[dvipsnames]{xcolor}       % colours; dvipsnames adds ~70 named colours
                                      % (NavyBlue, ForestGreen, ...)
\definecolor{linkblue}{RGB}{0,51,153} % dark blue used for links (see config.tex)
\usepackage{tikz}                     % drawing package; the cover uses it to place
                                      % the logos at exact positions on the page

% --- Typography ------------------------------------------------------------------------
\usepackage{microtype}    % small spacing adjustments for more even text and
                          % fewer bad line breaks: character protrusion (margin
                          % kerning) and, under pdfLaTeX, slight font expansion
\usepackage{xurl}         % lets long URLs (e.g. in references) break at any
                          % character instead of running into the margin; must
                          % load before hyperref
    % general packages (language, maths, graphics)
% ==================================================================================================
% preamble/fonts.tex: document fonts. Compile with XeLaTeX.
%
% Aptos is loaded from the font files in assets/font-aptos/
% Path: folder holding the files, relative to main.tex. Extension: added to each file name.
% Maths stays in the default Latin Modern maths font.
% ==================================================================================================

\usepackage{fontspec}
% Sets the body font
\setmainfont{Aptos}[
  Path           = assets/font-aptos/,
  Extension      = .ttf,
  UprightFont    = *,   % "*" stands for the base name "Aptos"
  BoldFont       = *-Bold,
  ItalicFont     = *-Italic,
  BoldItalicFont = *-Bold-Italic,
]
% Sets the sans-serif font to the same as the body font
\setsansfont{Aptos}[
  Path           = assets/font-aptos/,
  Extension      = .ttf,
  UprightFont    = *,
  BoldFont       = *-Bold,
  ItalicFont     = *-Italic,
  BoldItalicFont = *-Bold-Italic,
]
       % fonts
% ==================================================================================================
% preamble/layout.tex: page geometry, line spacing, headers/footers and heading styles.
% ==================================================================================================

% ======================================== page geometry ===========================================
\usepackage{geometry}
\geometry{
    a4paper,
    margin     = \thesismargin,
    headheight = 15pt,  % room for the running header
}

% ================================ line spacing and page breaking ==================================
\usepackage{setspace}   % line spacing; keeps footnotes and float contents single-spaced
\thesislinespacing      % body spacing
\raggedbottom           % let pages end at different heights
\widowpenalty=10000     % never leave the last line of a paragraph alone at the top of a page
\clubpenalty=10000      % never leave the first line of a paragraph alone at the bottom of a page
\setlength{\emergencystretch}{3em}  % stretch the spaces on lines without clean breaks

% ====================================== headers and footers =======================================
% fancyhdr defines "page styles": what goes in the header and footer.
%   \fancyhf{}             clear all header and footer fields
%   \fancyhead[L/C/R]{..}  left / centre / right of the header
%   \fancyfoot[L/C/R]{..}  left / centre / right of the footer
%   \headrulewidth         thickness of the line under the header (0pt = none)
%   \leftmark              holds the current chapter ("Chapter 3. Literature Review";
%   \rightmark             holds the current section
%   \thepage               prints the page number in the current style in the body font
\usepackage{fancyhdr}

% "thesis" style, used for normal pages in the main matter
\fancypagestyle{thesis}{%
    \fancyhf{}%
    \fancyhead[L]{\small\nouppercase{\leftmark}}% stops default ALL CAPS chapter headers
    \fancyfoot[C]{\thepage}%
    \renewcommand{\headrulewidth}{0.4pt}%
}

% "plain" style, used automatically on the first page of every chapter and on the front matter
\fancypagestyle{plain}{%
    \fancyhf{}%
    \fancyfoot[C]{\thepage}%
    \renewcommand{\headrulewidth}{0pt}%
}

\pagestyle{plain}                       % front matter: page numbers only
\appto\mainmatter{\pagestyle{thesis}}   % switch to the header style when \mainmatter runs

% ======================================== heading styles ==========================================
% titlesec controls how chapter and section titles look.
% \titleformat{command}[shape]{format}{label}{sep}{before-title}
%   shape           display = label on its own line above the title
%                   block   = label and title on one line
%   format          font applied to label and title
%   label           what prints as the number ("Chapter 1")
%   sep             space between label and title
%   before-title    code applied to the title text only (here: its size)
\usepackage{titlesec}

% Numbered chapters: "Chapter 1" (Large, bold) above the title (huge, bold).
% In the appendix \chaptertitlename becomes "Appendix", giving "Appendix A".
\titleformat{\chapter}[display]
    {\normalfont\bfseries}{\Large\chaptertitlename\ \thechapter}{10pt}{\huge}

% Unnumbered chapters (Acknowledgements, Abstract, Contents, lists,  References): title only.
\titleformat{name=\chapter,numberless}[block]
    {\normalfont\bfseries}{}{0pt}{\huge}

% \titlespacing*{command}{left indent}{space above}{space below}
% Default LaTeX puts a large gap above chapter titles; this removes it.
\titlespacing*{\chapter}{0pt}{0pt}{24pt}

% Simpler form for sections: \titleformat*{command}{format}.
% Sizes from large to small: \huge \LARGE \Large \large \normalsize.
\titleformat*{\section}{\normalfont\Large\bfseries}             % 1.1
\titleformat*{\subsection}{\normalfont\large\bfseries}          % 1.1.1
\titleformat*{\subsubsection}{\normalfont\normalsize\bfseries}  % unnumbered, see below

% ======================================== numbering depth =========================================
% Levels: 0 = chapter, 1 = section, 2 = subsection, 3 = subsubsection.
\setcounter{secnumdepth}{2} % number headings down to subsections (1.2.3)
                            % subsubsections get no number
\setcounter{tocdepth}{2}    % list headings down to subsections in the table of contents
      % margins, spacing, headers, heading styles
% =============================================================================
% preamble/floats.tex: figures, tables and algorithms ("floats": LaTeX moves
% them to where they fit best, so refer to them by number, not "below").
%
% Numbering is per chapter: Figure 3.1, Table 3.1, Algorithm 3.1. The book
% class does this for figures and tables; the [chapter] option below does it
% for algorithms.
%
% Conventions used here (and in most CS dissertations):
%   - table captions go ABOVE the table (put \caption before \begin{tabular})
%   - figure captions go BELOW the figure (put \caption after \includegraphics)
%   - put \label straight after \caption, so the label gets the right number
% Example code for all three is in README.md.
% =============================================================================

% --- Captions ------------------------------------------------------------------------------
\usepackage{caption}
\captionsetup{
  labelfont = bf,       % "Figure 3.1" in bold
  labelsep  = colon,    % "Figure 3.1: Caption text"
  font      = small,    % caption text one size smaller than body text
}
\captionsetup[table]{position=top}  % tells caption that table captions are
                                    % above, so it puts the gap below them

% --- Tables ---------------------------------------------------------------------------------
\usepackage{booktabs}   % professional horizontal rules: \toprule, \midrule,
                        % \bottomrule (avoid vertical lines in tables)
\usepackage{tabularx}   % tabularx environment: table of a set total width with
                        % X columns that wrap text and share the spare space
\usepackage{array}      % extra column options, e.g. p{3cm} with
                        % >{\raggedright\arraybackslash}

% --- Algorithms -----------------------------------------------------------------------------
% Two packages work together:
%   algorithm      the floating box with a caption and number, and the
%                  \listofalgorithms command used in frontmatter/lists.tex
%   algpseudocode  the pseudocode itself inside \begin{algorithmic}:
%                  \Require, \Ensure, \State, \If{..}..\EndIf, \For{..}..\EndFor,
%                  \While, \Function{Name}{args}..\EndFunction, \Return, \Comment
% \begin{algorithmic}[1] numbers every line; \label a line to \cref it.
\usepackage[chapter]{algorithm}  % [chapter]: number algorithms per chapter
\usepackage{algpseudocode}
      % captions, tables, algorithms
% =============================================================================
% preamble/references.tex: IEEE citations and reference list.
%
% HOW IT WORKS
% biblatex formats citations and the reference list; biber is the helper
% program that reads references.bib and sorts/numbers the entries (Overleaf
% runs it automatically). biblatex-ieee supplies the IEEE style.
% The reference list is printed by backmatter/references.tex.
%
% CITING IN THE TEXT
% citestyle = ieee-comp puts each number in its own brackets and collapses
% runs of consecutive numbers into a range:
%   \cite{a}                       [1]
%   \cite[p.~3]{b}                 [2, p. 3]      (\cite[3]{b} adds the "p." itself)
%   \cite[pp.~3--5]{b}             [2, pp. 3-5]
%   \cite[Sec.~4.2]{b}             [2, Sec. 4.2]
%   \cite{a,d,f}                   [1], [4], [6]
%   \cite{a,b,c,d}                 [1]-[4]        (consecutive numbers)
%   \cites[p.~3]{a}[Sec.~2]{b}     [1, p. 3], [2, Sec. 2]  (a note per source)
%   \textcite{b}                   Kumar et al. [2]  (author as sentence subject)
% "~" is a non-breaking space so "p." and "3" stay on one line.
% Every number is a link to its entry in the reference list (via hyperref).
% Numbers follow the order in which sources are first cited, as IEEE requires.
% (The plain "ieee" citestyle would print [1], [2], [3], [4] with no range.)
% =============================================================================

\usepackage[
  backend   = biber,      % biber (not the older bibtex) processes references.bib
  style     = ieee,       % IEEE reference-list format (initials, "vol.", "pp.",
                          % "doi:"); entries listed in citation order
  citestyle = ieee-comp,  % in-text [n] with ranges, as described above
  doi       = true,       % print DOIs in the reference list
  url       = true,       % print URLs (for sources without a DOI)
  isbn      = false,      % don't print ISBNs (IEEE references normally omit them)
]{biblatex}

\addbibresource{references.bib}  % the database file (name includes .bib)
  % IEEE citations with biblatex
% =============================================================================
% preamble/links.tex: clickable links and named cross-references.
%
% hyperref makes every citation, cross-reference, table-of-contents entry,
% glossary term and URL clickable, and builds the PDF bookmarks panel. It
% changes many other packages' commands, so it loads after almost everything;
% only cleveref (right below) and glossaries-extra (glossary.tex) come later.
% =============================================================================

% --- hyperref ----------------------------------------------------------------------
\usepackage[
  colorlinks,          % colour link text instead of drawing boxes around it
  bookmarksnumbered,   % PDF bookmarks show chapter/section numbers
]{hyperref}

\hypersetup{
  linkcolor  = \thesislinkcolour,  % internal links: ToC, \cref, glossary
  citecolor  = \thesislinkcolour,  % citation numbers [1]
  urlcolor   = \thesislinkcolour,  % web addresses
  pdftitle   = {\thesistitle},     % PDF metadata, shown in the viewer's title bar
  pdfauthor  = {\thesisauthor},    % and in File -> Properties
  pdfsubject = {\thesisdegree{} dissertation},
}

% --- cleveref ------------------------------------------------------------------------
% Writes the reference type for you, so numbers and names can't get out of
% step. Label something, then refer to it:
%   \label{fig:pipeline}  ->  \cref{fig:pipeline}   "Figure 3.1"
%                             \Cref{fig:pipeline}   same, for sentence starts
%   \cref{fig:a,fig:b}        "Figures 3.1 and 3.2"
%   \cref{eq:rev2}            "Equation (2.1)"
% Label prefixes used in this template, so labels stay readable:
%   ch: chapter   sec: section   fig: figure   tab: table
%   alg: algorithm   eq: equation   app: appendix
% Plain \ref{label} still works and prints just the number ("3.1").
\usepackage[
  capitalise,   % "Figure 3.1" rather than "figure 3.1" mid-sentence
  noabbrev,     % "Figure", "Equation" rather than "Fig.", "Eq."
]{cleveref}

% Names cleveref prints for types it doesn't name the way we want:
% \crefname{type}{singular}{plural}
\crefname{algorithm}{Algorithm}{Algorithms}  % algorithm floats
\crefname{ALG@line}{line}{lines}             % labelled lines inside algorithmic
\crefname{appendix}{Appendix}{Appendices}    % chapters after \appendix
       % hyperref and cleveref; keep after the files above
% =============================================================================
% preamble/glossary.tex: abbreviations and glossary (glossaries-extra).
%
% WHAT YOU GET
% Appendix A has two numbered sections, printed by
% backmatter/appendix-glossary.tex:
%   A.1 Abbreviations  short form -> long form   (e.g. MCP  Model Context Protocol)
%   A.2 Glossary       term -> definition
% Entries are defined in glossary/entries.tex. Only entries used somewhere
% in the text are printed, alphabetically. Sorting is done by LaTeX itself
% (\makenoidxglossaries), so no extra program is needed on Overleaf.
%
% USING ENTRIES IN THE TEXT (every use is a link to the appendix entry)
%   \gls{mcp}          first use "Model Context Protocol (MCP)", later "MCP"
%   \glspl{llm}        plural: first use "large language models (LLMs)"
%   \Gls{...} \Glspl   capitalised first letter, for the start of a sentence
%   \glsxtrshort{mcp}  always "MCP"
%   \glsxtrlong{mcp}   always "Model Context Protocol"
%   \glsxtrfull{mcp}   always "Model Context Protocol (MCP)"
%   \gls{def-mcp}      "Model Context Protocol", linking to its definition
%   \gls{rug-pull}     a glossary term, linking to its definition
%   \glsfmtshort{mcp}  "MCP" in a chapter/section title or caption; there, \gls
%                      would use up the first-use expansion and put links in
%                      the table of contents
%   \glsreset{mcp}     make the next \gls{mcp} expand again (e.g. per chapter)
%
% ABBREVIATION -> DEFINITION LINK
% For an abbreviation that also needs a definition, \newtermabbr (defined
% below) creates two linked entries. Its line in A.1 ends with
% "(see definition)", which links to the definition in A.2.
% =============================================================================

% --- Package options ---------------------------------------------------------------------
\usepackage[
  abbreviations,               % create the separate "abbreviations" list (A.1)
  nonumberlist,                % don't list page numbers after each entry
  section    = section,        % print each list as a \section, so both sit
                               % inside the appendix chapter
  numberedsection = autolabel, % number those sections (A.1, A.2) and give each
                               % a \label: \glsautoprefix + list name
  toc        = false,          % numbered sections reach the table of contents
                               % anyway; this avoids a duplicate entry
  nogroupskip,                 % no extra gap between letter groups (A, B, ...)
  stylemods,                   % load glossaries-extra's fixes to the list styles
                               % (needed for the "(see definition)" hook below)
]{glossaries-extra}

\renewcommand*{\glsautoprefix}{app:}  % section labels become app:abbreviations
                                      % (A.1) and app:main (A.2), for \cref
\makenoidxglossaries                  % sort and print lists without makeindex

% --- Appearance ----------------------------------------------------------------------------
% long-short: first use prints "long form (SHORT)", later uses "SHORT".
% Other styles: short-long ("SHORT (long form)"), long-short-sc (small caps).
\setabbreviationstyle{long-short}             % ordinary abbreviations
\setabbreviationstyle[termabbr]{long-short}   % abbreviations made by \newtermabbr
                                              % (their own category, see below)
\renewcommand*{\glsnamefont}[1]{\textbf{#1}}  % entry names in bold in both lists
\setlength{\glsdescwidth}{0.75\linewidth}     % width of the long-form column in
                                              % A.1 (the "long" list style)

% --- \newtermabbr: abbreviation + definition, linked -----------------------------------------
% Usage: \newtermabbr{label}{SHORT}{long form}{definition}
% Creates:
%   label      an abbreviation (A.1). Its "category" is set to termabbr so the
%              two hooks below apply only to these entries.
%   def-label  a glossary term (A.2) named "long form (SHORT)" with the
%              definition; in the text \gls{def-label} prints just "long form".
% #1..#4 are the four arguments in order.
\newcommand{\newtermabbr}[4]{%
  \newglossaryentry{def-#1}{name={#3 (#2)}, text={#3}, description={#4}}%
  \newabbreviation[category=termabbr]{#1}{#2}{#3}%
}

% Helper commands for the hooks: add the definition entry to the Glossary,
% and print the "(see definition)" link to it.
\newcommand*{\addtermdefinition}[1]{\glsadd{def-#1}}
\newcommand*{\linktermdefinition}[1]{%
  \space(\emph{see} \glshyperlink[definition]{def-#1})}

% Hook 1, runs after every \gls-type use of a termabbr abbreviation: also
% mark its definition as used, so the definition always appears in A.2 and
% the "(see definition)" link always has a target.
% \glslabel holds the label just used; \expandafter turns it into its text
% (e.g. "mcp") before the helper builds "def-mcp".
\glsdefpostlink{termabbr}{\expandafter\addtermdefinition\expandafter{\glslabel}}

% Hook 2, runs after the description of each termabbr entry in the printed
% list: append "(see definition)". \glscurrententrylabel is the label of the
% entry being printed.
\glsdefpostdesc{termabbr}{%
  \expandafter\linktermdefinition\expandafter{\glscurrententrylabel}}

% --- Entries -------------------------------------------------------------------------------
% Must come after the settings above.
% =============================================================================
% glossary/entries.tex: abbreviation and glossary entries.
% Loaded at the end of preamble/glossary.tex (see that file for all the commands you can use in the text).
%
% Only entries used in the text appear in Appendix A, sorted alphabetically, so defining an entry you never use is harmless. Keep this file grouped by kind and alphabetical by label so entries are easy to find.
%
% THE THREE KINDS OF ENTRY
%
% 1. Abbreviation only -> A.1 Abbreviations
%      \newabbreviation{label}{SHORT}{long form}
%    Write the long form in lower case unless it is a proper noun; \Gls capitalises it at the start of a sentence.
%      \newabbreviation[longplural={...}, shortplural={...}]{label}{SHORT}{long form}
%
% 2. Defined term only -> A.2 Glossary
%      \newglossaryentry{label}{name={term}, description={definition}}
%    Optional keys: plural={...} for an irregular plural; sort={...} to sort under a different spelling.
%
% 3. Abbreviation that also has a definition -> both lists, linked
%      \newtermabbr{label}{SHORT}{long form}{definition}
%    In the text, \gls{label} gives the abbreviation and \gls{def-label} the defined term. The A.1 line ends "(see definition)" and links to A.2.
%
% Labels: lower case, no spaces; hyphens are fine (e.g. rug-pull). The label is only used in \gls{...}, never printed. Commas or "=" inside a value need braces: description={a, b, and c}.
% =============================================================================

% --- 1. Abbreviations ----------------------------------------------------------------------


% --- 3. Abbreviations with a definition ------------------------------------------------------
% The fourth argument (the definition) may start on the next line, but there must be no blank line between it and the third argument.


% --- 2. Defined terms ------------------------------------------------------------------------


    % glossaries-extra must load after hyperref

\begin{document}

% ========================================= front matter ===========================================
% The cover gets letter page numbers (A) that are never printed. This keeps its internal page ID
% different from page i and page 1, so hyperref does not warn about duplicate page destinations.
\pagenumbering{Alph}
% =============================================================================
% frontmatter/cover.tex: cover page.
%
% Follows Warwick's "Student Thesis Cover - Instructions for LaTeX"
% (_reference/Instructions for LaTeX.pdf):
%   - crest at the top, wordmark at the bottom, everything centred
%   - title: Aptos 20pt bold; "BY": 12pt capitals; name: 20pt; other text: 14pt
% The cover uses the document font set in preamble/fonts.tex (Aptos where
% the compiler supports it). The text comes from config.tex; edit it there,
% not here.
%
% HOW IT WORKS
% A TikZ picture with the "overlay" option draws on the page without taking
% up space, and "remember picture" lets it use (current page.north) and
% (current page.south), the centres of the top and bottom page edges, as
% anchor points. Each \node is placed relative to one of those, so positions
% are fixed on the paper regardless of margins. The distances were measured
% from the example cover in the instructions PDF. Overlay positions need two
% compile runs; Overleaf does this automatically.
%
% TO ADJUST
%   move an item:   change its yshift (negative = further down from the top)
%   resize a logo:  change width= in its \includegraphics
%   title width:    text width=0.62\paperwidth (narrower = more lines)
% =============================================================================

\begin{titlepage}  % a page with no header, footer or page number
\setstretch{1}     % single line spacing on the cover (body text uses 1.5)
\begin{tikzpicture}[remember picture, overlay]

  % Crest: top edge 20 mm below the top of the paper, 35 mm wide.
  % anchor=north: the node's top-centre sits at the given point.
  \node[anchor=north] at ([yshift=-20mm]current page.north)
    {\includegraphics[width=35mm]{warwick-coat-of-arms}};  % assets/ folder

  % Text block: top edge 120 mm below the top of the paper. A longer title
  % pushes the rest of the block down without moving the logos.
  % text width makes the node a fixed-width paragraph box, so \par, \vspace
  % and automatic line wrapping work inside it.
  \node[anchor=north, text width=0.62\paperwidth]
    at ([yshift=-120mm]current page.north) {%
      \centering
      % \fontsize{size}{line distance}\selectfont sets an exact point size.
      {\fontsize{20}{25}\selectfont\bfseries \covertitle\par}  % title, 20pt bold
      \vspace{8pt}
      {\fontsize{12}{15}\selectfont BY\par}                    % "BY", 12pt caps
      \vspace{6pt}
      {\fontsize{20}{25}\selectfont \thesisauthor\par}         % name, 20pt
      \vspace{24pt}
      {\fontsize{14}{19}\selectfont                            % remaining text, 14pt
        A dissertation submitted in partial fulfilment of the\\
        requirements for the degree of\par
        \vspace{6pt}
        \thesisdegree\par
        \thesisdepartment\par
        \vspace{12pt}
        Supervisor: \thesissupervisor\par
        \vspace{12pt}
        \thesisdate\par}
    };

  % Wordmark: bottom edge 18 mm above the bottom of the paper, 65 mm wide.
  % anchor=south: the node's bottom-centre sits at the given point.
  \node[anchor=south] at ([yshift=18mm]current page.south)
    {\includegraphics[width=65mm]{warwick-wordmark}};      % assets/ folder

\end{tikzpicture}
\end{titlepage}


% Page numbers restart at i, ii ... and chapters are unnumbered
\frontmatter
% ==================================================================================================
% frontmatter/acknowledgements.tex
% ==================================================================================================

\chapter*{Acknowledgements}
\pdfbookmark[0]{Acknowledgements}{acknowledgements}

[Acknowledgements to be written.]

% ==================================================================================================
% frontmatter/abstract.tex
% Unnumbered chapter with a PDF bookmark but no table-of-contents entry
% \chapter* = unnumbered chapter: new page, big title, no number and no table-of-contents entry.
% \pdfbookmark[level]{text}{id} still adds it to the PDF viewer's bookmarks panel; level 0 = same
% level as chapters, id = any unique name.
% ==================================================================================================

\chapter*{Abstract}
\pdfbookmark[0]{Abstract}{abstract}

[Abstract to be written.]

% Optional keywords line:
% \medskip                          % small vertical gap
% \noindent\textbf{Keywords:} ...   % \noindent: no paragraph indent

% =============================================================================
% frontmatter/lists.tex: table of contents, then the lists of tables, figures and algorithms.
%
% The lists are generated from \caption commands (tables, figures, algorithms) and headings (contents), via auxiliary files written on the previous compile run, so new items appear after the next run.
% =============================================================================

\begingroup          % start a group: the spacing change below stays inside it
\setstretch{1}       % lists are single-spaced even though body text is 1.5

\tableofcontents     % "Contents"; not listed in itself

% Each block below:
%   \clearpage            finish the previous list's page
%   \phantomsection       invisible anchor so the ToC link jumps to this page
%   \addcontentsline{toc}{chapter}{title}  add a chapter-level ToC entry
%   \listof...            print the list (starts with an unnumbered chapter
%                         heading on the fresh page)
% The ToC entry is added before the list so it records the right page.

\clearpage\phantomsection
\addcontentsline{toc}{chapter}{\listtablename}      % "List of Tables"
\listoftables

\clearpage\phantomsection
\addcontentsline{toc}{chapter}{\listfigurename}     % "List of Figures"
\listoffigures

\clearpage\phantomsection
\addcontentsline{toc}{chapter}{\listalgorithmname}  % "List of Algorithms"
\listofalgorithms                                   % from the algorithm package

\endgroup            % back to the body line spacing
 % table of contents, lists of tables/figures/algorithms

% ========================================= main matter ============================================
% Page numbers restart at 1 in arabic numerals, chapters are numbered, and the running header is
% switched on (see preamble/layout.tex).
\mainmatter
% =============================================================================
% chapters/ch1-introduction.tex: Chapter 1.
%
% HOW A CHAPTER FILE WORKS (applies to every file in chapters/)
% - The file is pulled in by \include in main.tex, which starts a new page.
% - \chapter{Title} prints "Chapter 1" + title and starts the running header; \label{ch:...} right after it lets you write \cref{ch:introduction} ("Chapter 1") anywhere.
% - \section / \subsection are numbered (1.1, 1.1.1) and listed in the table of contents; \subsubsection is unnumbered. Add \label{sec:...} after each one you want to refer to.
% - A heading that needs an abbreviation: \section{Why \glsfmtshort{mcp} ...} (plain "MCP" also works but won't count as a use of the glossary entry).
% - Starred versions (\section*{...}) are unnumbered and not in the contents.
% - Citations: \cite{kumar2018rev2}, \cite[p.~3]{key} (keys: references.bib).
%   Glossary: \gls{mcp} (entries: glossary/entries.tex).
%   Figures, tables, algorithms: see README.md.
% - A blank line starts a new paragraph; a single line break does not.
% - % starts a comment; to print a percent sign write \%.
%   Other characters that need a backslash: \& \# \_ \$ \{ \}.
% =============================================================================

\chapter{Introduction}\label{ch:introduction}

% =============================================================================
% chapters/ch2-background.tex: Chapter 2.
% =============================================================================

\chapter{Background}\label{ch:background}

% TEMPORARY placeholder content: the minimum needed to clear the empty-bibliography, empty-glossary and unused-table-caption warnings. Replace when writing the chapter.
% \gls{mcp}          uses the MCP abbreviation (fills Abbreviations) and, through \newtermabbr, its definition (fills Glossary)
% \cite{...}         cites one source (fills References)
% table + \caption   uses the table caption setup from preamble/floats.tex
Placeholder text using \gls{mcp} \cite{hou2026model}.

\begin{table}[tb]
  \centering
  \caption{Placeholder table.}\label{tab:placeholder}
  \begin{tabular}{l}
    \toprule
    Placeholder \\
    \bottomrule
  \end{tabular}
\end{table}

% =============================================================================
% chapters/ch3-literature-review.tex: Chapter 3.
% Section headings follow drafts/outline-part1-intro-background-litreview.md.
% How chapter files work (headings, labels, citations, glossary): see the
% comment block at the top of chapters/ch1-introduction.tex.
% =============================================================================

\chapter{Literature Review}\label{ch:literature-review}

\section{Malicious MCP servers and the limits of current defences}\label{sec:lr-malicious-mcp}

\section{Attacks that exploit reputation history}\label{sec:lr-history-attacks}

\section{Graph-based rating-fraud detection, and why REV2}\label{sec:lr-graph-fraud}

\section{Approaches to behaviour change in trust assessment}\label{sec:lr-behaviour-change}

\section{Time windows on a coupled graph score: constraints and remaining choices}\label{sec:lr-windows}

\section{Computing and thresholding the delta}\label{sec:lr-delta}

\section{The one-directional reporting rule}\label{sec:lr-reporting-rule}

\section{Evaluating trust mechanisms under simulated attacks}\label{sec:lr-evaluation}

\section{Synthesis, research gap and design decisions}\label{sec:lr-synthesis}

% =============================================================================
% chapters/ch4-methodology.tex: Chapter 4.
% Sections to be added once the methodology outline is settled.
% How chapter files work (headings, labels, citations, glossary): see the
% comment block at the top of chapters/ch1-introduction.tex.
% =============================================================================

\chapter{Methodology}\label{ch:methodology}


% ========================================== references ============================================
% Printed after the last chapter and listed in the table of contents.
% =============================================================================
% backmatter/references.tex: the reference list.
%
% \printbibliography prints every source cited with \cite and friends, in
% IEEE format and in order of first citation (settings in
% preamble/references.tex; entries in references.bib). Uncited entries in
% references.bib are not printed. If nothing is cited yet, nothing prints
% (not even the heading).
%   heading = bibintoc  unnumbered chapter heading, listed in the table of
%                       contents
%   title = {...}       heading text
% =============================================================================

\printbibliography[heading=bibintoc, title={References}]


% ========================================== appendices ============================================
% Chapters from here on are lettered (Appendix A, B, ...) instead of numbered.
\appendix
% =============================================================================
% backmatter/appendix-glossary.tex: Appendix A, Glossary and Abbreviations.
%
% Because this file comes after \appendix in main.tex, \chapter prints
% "Appendix A". Each \printnoidxglossary below becomes a numbered section
% (A.1, A.2) holding one list, generated from glossary/entries.tex (settings
% in preamble/glossary.tex). Only entries used in the text are printed; a
% list with no used entries prints nothing.
%
% Referencing from the text:
%   \cref{app:glossary}       "Appendix A"
%   \cref{app:abbreviations}  "Appendix A.1"
%   \cref{app:main}           "Appendix A.2"
%
% Options of \printnoidxglossary:
%   type   which list: \glsxtrabbrvtype = abbreviations, main = glossary terms
%   title  section heading
%   style  layout: long    = two-column table (short form | long form)
%                  altlist = name on its own line, definition indented below
%          Other styles: https://ctan.org/pkg/glossaries (glossary style list)
% =============================================================================

\chapter{Glossary and Abbreviations}\label{app:glossary}

\printnoidxglossary[type=\glsxtrabbrvtype, title={Abbreviations}, style=long]  % A.1
\printnoidxglossary[type=main, title={Glossary}, style=altlist]                % A.2


\end{document}
